\documentclass{article}
\usepackage{graphicx}
\usepackage{float}
\setlength{\parindent}{0pt}
\usepackage{caption}
\usepackage{fancyvrb}
\usepackage{tcolorbox}
\usepackage[a4paper, total={6in, 8.5in}]{geometry}
\setlength{\footskip}{80pt}
\usepackage{tabularx}
\usepackage{booktabs}
\usepackage{enumitem}

\title{Project - Network Redesign}
\author{Aleksander Kłap}
\date{08.06.2026}

\begin{document}

\maketitle
\pagebreak

\section{Logical Network Redesign}

TODO: IN PROGRESS
TODO: DIAGRAM HERE

%%\begin{figure}[H]
%%    \centering
%%    \includegraphics[width=1\textwidth]{img/conduit_diagram.pdf}
%%    \caption{Zone \& Conduit Diagram}
%%\end{figure}

The current flat /18 network is replaced with a layered zone architecture 
based on the IEC 62443 zone and conduit model. 
Each zone groups devices with similar operational functions. 
All inter-zone traffic passes through a firewall and is restricted to explicitly whitelisted traffic paths.

\subsection{Logic Zones}
\subsubsection{IT Zone}
The IT Zone contains nodes handling business logic and order management. It have no direct
routed path to any OT device - all IT to OT traffic goes through the IDMZ.
\subsubsection{IDMZ}
The IDMZ acts as a buffer between IT and OT devices. No connecection is allowed 
to pass from the IT to OT and vice versa. Every traffic exchange terminates inside IDMZ,
and its re-innitiated to the other side. This solution implements IEC 62443-3-3 SR 5.1 principle.
\subsubsection{OT Cell Zone}
The OT Cell Zone Contains devices supporting production, but not beeing part of the physical production
process. The Enigneering Workstation is placed in this logical zone, however, in physical
implementation, this node will be further isolated into separate VLAN.
\subsubsection{OT Production Zone}
The OT Production Zone contains devices beeing part of the direct physical production process.
Only inbound traffic allowed in this zone are reads and writes from the proxy services, and 
S7comm+ from the engineering workstation.

\subsection{New Nodes}
\subsubsection{OPC UA proxy}
In the AS-IS architecture, MES can directly connect to the PLC's OPC UA servers over the flat network.
OPC UA proxy node aggregates all PLC's OPC UA servers traffic into single endpoint inside IDMZ.
The MES connects only to proxy, that is also responsible for enforcing rules on 
what nodes are readable and writable from the MES. Furthermore, it logs all sessions.
This solution satisfies IEC 62443-3-3 SR 3.2.
\subsubsection{Nginx}
In AS-IS architecture, MES can directly access unencrypted HTTP services from the OT nodes.
As stated in security requirements, all not strictly required HTTP servers of OT devices will be 
disabled. However, for required HTTP traffic (e.g PLC's alarm table <- MES) the Nginx node is
introduced in IDMZ. MES connects to the Nginx over HTTPS TLS 1.2, terminates the connection, access
the requested HTTP node server and forwards it to the MES. Nginx enforces authentication,
logging all requests and block requests for services that should not be reached from IT zone.
\subsubsection{Dedicated MQTT broker}
In the AS-IS architecture MQTT CVX device publish messages over a flat network to the MES.
MQTT broker node placed in IDMZ, aggregates MQTT traffic from the CVX device using TLS 1.2 encryption,
and it restrict inbound traffic to publish messages only. MES is subscribes to the
MQTT broker, preventing direct communication between OT node and IT node. This solution
satisfies IEC 62443-3-3 SR 1.1 and SR 3.1.
\subsubsection{Engineering Workstation}
In the AS-IS architecture, TIA Portal is run on the same PC as MES service - part of IT and reachable over flat network.
A dedicated Engineering Workstation is placed in the OT Cell Zone (physically on separate VLAN).
This workstation purpose is to program and configure PLC's and HMI's with separation
from IT nodes. Traffic is restricted only to TCP 102 on PLC's and HMI's. This enforces
least-privilage and IEC 62443-2-1 §4.3.3.
\subsubsection{Syslog}
In the AS-IS architecture, there is no centralized logging system implemented.
Syslog node collects all logs from the firewalls, managed switches, OPC UA and HTTP proxy, and MQTT broker.
This provides post-attack analysis and satisfies IEC 62443-3-3 SR 6.1 and NIS2 Article 21(2)(f). 
\subsubsection{Backup}
In AS-IS architecture, there is no documented backup procedure. The backup node recieves
configuration backups from the Engineering Workstation after every change and daily database
backups from the MES. It is placed on separate VLAN inside IT zone, separated from
indbound traffic from OT. Satisfies IEC 62443-2-1 §4.3.4 and ISO/IEC 27001 A.12.3.
\subsubsection{Log \& Backup proxy}



\subsection{Firewalls}
This section only describres briefly the high-level purpose of the firewall nodes.
Further specification on permitted conduits is described in TODO
\subsubsection{Edge Firewall}
Terminates all VPN tunnels from remote engineers and vendor support, 
enforces MFA before any traffic enters the IT zone, and blocks all unsolicited inbound 
connections from the internet. Default-deny on all inbound. 
Satisfies IEC 62443-2-1 §4.3.3.6 and NIS2 Article 21(2)(i).
\subsubsection{IT - DMZ Firewall}
Controls all traffic between the IT zone and the IDMZ. 
\subsubsection{DMZ - OT Cell Firewall}
Controls traffic between the IDMZ and the OT cell zone. 
\subsubsection{OT Cell - OT Production Firewall}
Controls traffic between the OT cell zone and the six station VLANs. 
OPC UA proxy connections arrive from the IDMZ side, not the cell zone side, 
so they do not pass this firewall. Profinet never crosses this boundary — 
it remains Layer 2 within each station. All other traffic denied. 
Satisfies IEC 62443-3-3 SR 5.1 and the least privilege principle.

\section{Network Implementation}
TODO: DIAGRAM HERE

\subsection{VLAN Segmentation}
The redesigned factory network is implemented using VLAN-based segmentation 
on managed industrial switches. Each security zone is assigned a dedicated 
VLAN and IP subnet. Inter-VLAN communication is not performed by the switch 
itself but is routed through a central industrial firewall.
The firewall acts as the default gateway for all VLANs and enforces security 
policies between zones. A VLAN trunk connection between the firewall and the 
managed switch transports traffic for all security zones while preserving 
logical separation through IEEE 802.1Q VLAN tagging. \\

\renewcommand{\arraystretch}{1.3} % 30% more row height
\begin{tabularx}{\textwidth}{>{\bfseries}p{2.5cm} X X}
    VLAN & \textbf{Zone} & \textbf{Nodes} \vspace{0.5em}\\
\toprule
    10 & IT & MES PC, Data Storage PC \\
    20 & IDMZ & Proxy PC\\
    30 & OT Cell & COBOL PC, MQTT CVX \\
    40 & OT Cell - Engineering & Engineering Workstation  \\
    101 & OT Prod - Station 1 & PLC1, HMI1, UR3e Robot Arm Controller, TURCK position sensor \\
    102 & OT Prod - Station 2 & PLC2, HMI2, I/O device, Festo Controller \\
    103 & OT Prod - Station 3 & PLC3, HMI3, SensoPart Camera \\
    104 & OT Prod - Station 4 & PLC4, HMI4 \\
    105 & OT Prod - Station 5 & PLC5, HMI5 \\
    106 & OT Prod - Station 6 & PLC6, HMI6, Festo controller \\
\end{tabularx}

\subsection{Restricted Routing}
\end{document}

